\documentclass[12pt,a4paper]{article}

% ---------------------------------------------------------------
% Documentação do Sistema e da API — Manus Hospital
% Compilar no Overleaf com pdfLaTeX (padrão).
% ---------------------------------------------------------------
\usepackage[utf8]{inputenc}
\usepackage[T1]{fontenc}
\usepackage[brazil]{babel}
\usepackage{lmodern}
\usepackage{microtype}
\sloppy
\setlength{\parindent}{0pt}
\setlength{\parskip}{6pt}
\usepackage[a4paper,margin=2.5cm]{geometry}
\usepackage{xcolor}
\usepackage{listings}
\usepackage{booktabs}
\usepackage{tabularx}
\usepackage{longtable}
\usepackage{array}
\usepackage{etoolbox}
\AtBeginEnvironment{tabularx}{\small}
\AtBeginEnvironment{longtable}{\small}
\usepackage{enumitem}
\usepackage{titlesec}
\usepackage{fancyhdr}
\usepackage{tikz}
\usetikzlibrary{arrows.meta,positioning}
\usepackage[hidelinks]{hyperref}

\definecolor{azul}{HTML}{14213D}
\definecolor{verde}{HTML}{1F7A6F}
\definecolor{fundo}{HTML}{F4F6F8}
\definecolor{borda}{HTML}{D5DCE3}

\titleformat{\section}{\Large\bfseries\color{azul}}{\thesection}{0.6em}{}
\titleformat{\subsection}{\large\bfseries\color{verde}}{\thesubsection}{0.6em}{}

\pagestyle{fancy}
\fancyhf{}
\fancyhead[L]{\small Manus Hospital — Documentação do Sistema e da API}
\fancyhead[R]{\small\thepage}
\renewcommand{\headrulewidth}{0.4pt}

\newcolumntype{L}{>{\raggedright\arraybackslash}X}
\renewcommand{\arraystretch}{1.25}

\lstset{
  basicstyle=\ttfamily\small,
  backgroundcolor=\color{fundo},
  frame=single,
  rulecolor=\color{borda},
  breaklines=true,
  columns=fullflexible,
  keepspaces=true,
  showstringspaces=false,
  xleftmargin=4pt, xrightmargin=4pt,
  literate=
    {á}{{\'a}}1 {à}{{\`a}}1 {ã}{{\~a}}1 {â}{{\^a}}1
    {é}{{\'e}}1 {ê}{{\^e}}1 {í}{{\'i}}1
    {ó}{{\'o}}1 {õ}{{\~o}}1 {ô}{{\^o}}1 {ú}{{\'u}}1
    {ç}{{\c{c}}}1 {Ç}{{\c{C}}}1 {Á}{{\'A}}1 {É}{{\'E}}1
    {Ó}{{\'O}}1 {Ã}{{\~A}}1 {→}{{$\rightarrow$}}1
}
\newcommand{\cod}[1]{\texttt{#1}}

\begin{document}

% ======================= CAPA =======================
\begin{titlepage}
  \centering
  {\large INSTITUTO FEDERAL DO SUDESTE DE MINAS GERAIS\\ Campus Muriaé\par}
  \vspace{0.5cm}
  {\large Gestão da Tecnologia da Informação\\ Engenharia de Software\par}
  \vfill
  {\Huge\bfseries\color{azul} Manus Hospital\par}
  \vspace{0.4cm}
  {\LARGE Documentação do Sistema e da API\par}
  \vspace{0.4cm}
  {\large API REST em C\# / ASP.NET Core integrada ao aplicativo PHP\par}
  \vfill
  \begin{tabular}{ll}
    \textbf{Integrantes:} & Karla Pardócimo \\
                          & Amanda Garcia \\
                          & Gabriel Pimentel \\
  \end{tabular}
  \vfill
  {\large Muriaé — MG\\ 2026\par}
\end{titlepage}

\tableofcontents
\newpage

% ======================= 1 =======================
\section{Introdução}

Este documento descreve o \textbf{Manus Hospital}, sistema de gestão hospitalar desenvolvido na
disciplina anterior em PHP e MySQL e evoluído nesta etapa para uma arquitetura
\textbf{cliente/servidor}. Na versão anterior o aplicativo PHP acessava o banco de dados
diretamente. Agora uma \textbf{API REST em C\# / ASP.NET Core} concentra as regras de negócio,
o acesso aos dados, a autenticação JWT e o cache, e o aplicativo PHP passa a ser apenas um
cliente dessa API.

O documento cobre os itens exigidos na seção 9 da especificação do trabalho: descrição do
sistema, tecnologias, arquitetura, execução da API e do aplicativo, banco de dados,
autenticação JWT, estratégia de cache, endpoints, exemplos de requisições e respostas e
reprodução do projeto. Inclui também as principais decisões arquiteturais.

O código-fonte completo, o \cod{README.md} e os documentos complementares
(\cod{docs/REGRAS\_DE\_NEGOCIO.md} e \cod{docs/DECISOES\_ARQUITETURAIS.md}) estão no
repositório Git do projeto.

% ======================= 2 =======================
\section{Descrição do sistema}

O sistema cobre o fluxo de um hospital, do cadastro do paciente ao faturamento:

\begin{itemize}[itemsep=2pt]
  \item \textbf{Pacientes}, com validação de CPF e registro de alergias, verificadas em toda prescrição;
  \item \textbf{Consultas e fila de triagem}: agendamento, check-in, cancelamento e falta;
  \item \textbf{Triagem pelo Protocolo de Manchester}: a API calcula a cor de risco a partir dos
        sinais vitais e ordena a fila de atendimento por gravidade;
  \item \textbf{Prontuário eletrônico}: evolução clínica, prescrições com bloqueio por alergia e alerta
        de interação medicamentosa, exames e resultados;
  \item \textbf{Leitos e internações}: mapa de leitos, internação, transferência e alta;
  \item \textbf{Estoque de insumos}: entradas, saídas, estoque crítico e uso em paciente internado;
  \item \textbf{Faturamento}: consolidação de gastos, emissão de guias, auditoria de valores altos,
        pagamento e cancelamento;
  \item \textbf{Usuários e perfis} (Administrador, Médico, Enfermagem e Recepção), cada um com
        permissões diferentes.
\end{itemize}

Cada regra de negócio recebe um código (por exemplo \cod{RN-PAC-02} ou \cod{RN-TRI-05}), que
aparece nas mensagens de erro da API. A lista completa está em
\cod{docs/REGRAS\_DE\_NEGOCIO.md}.

% ======================= 3 =======================
\section{Tecnologias utilizadas}

\begin{tabularx}{\textwidth}{@{}lL@{}}
\toprule
\textbf{Parte} & \textbf{Tecnologia} \\
\midrule
Linguagem da API     & C\# no \textbf{.NET 10 (LTS)}, ASP.NET Core Web API com controllers \\
Acesso a dados       & ADO.NET com \textbf{MySqlConnector 3.0} e SQL parametrizado \\
Banco de dados       & \textbf{MariaDB 10.11} (Docker) ou MySQL 8 (XAMPP) \\
Autenticação         & \textbf{JWT} assinado com HMAC-SHA256 (\cod{JwtBearer}) \\
Senhas               & \textbf{PBKDF2-SHA256}, 100.000 iterações e salt por usuário \\
Cache                & \textbf{IMemoryCache}, estratégia \textit{cache-aside} \\
Documentação da API  & Swagger / OpenAPI (Swashbuckle) \\
Aplicativo cliente   & PHP 8.3 com cURL, Bootstrap 5 \\
Implantação          & \textbf{Docker} e Docker Compose \\
Versionamento        & Git \\
\bottomrule
\end{tabularx}

% ======================= 4 =======================
\section{Estrutura da arquitetura}

\subsection{Visão geral}

A Figura~\ref{fig:geral} mostra a arquitetura final, que segue o modelo da especificação: o
aplicativo cliente conversa com a API por HTTP, e somente a API acessa o banco de dados.

\begin{figure}[h]
\centering
\begin{tikzpicture}[
  caixa/.style={draw=azul, rounded corners, minimum width=2.9cm, minimum height=1.3cm,
                align=center, font=\small, fill=white},
  dest/.style={caixa, fill=verde!12, draw=verde},
  seta/.style={-{Stealth[length=3mm]}, thick, color=azul}]
  \node[caixa] (php) {Aplicativo PHP\\\scriptsize container \cod{web} :8080};
  \node[dest, right=2.4cm of php] (api) {API ASP.NET Core\\\scriptsize container \cod{api} :5080};
  \node[caixa, right=2.4cm of api] (db) {MariaDB\\\scriptsize container \cod{db} :3307};
  \draw[seta] (php) -- node[above, font=\scriptsize]{HTTP + JWT} (api);
  \draw[seta] (api) -- node[above, font=\scriptsize]{ADO.NET} (db);
\end{tikzpicture}
\caption{Aplicativo cliente, API e banco de dados em três containers.}
\label{fig:geral}
\end{figure}

\subsection{Clean Architecture}

A API está dividida em quatro projetos. As dependências apontam para dentro: o domínio não
conhece banco de dados, HTTP nem bibliotecas externas.

\begin{figure}[h]
\centering
\begin{tikzpicture}[
  caixa/.style={draw=azul, rounded corners, minimum width=3.4cm, minimum height=1cm,
                align=center, font=\small, fill=white},
  seta/.style={-{Stealth[length=3mm]}, thick, color=azul}]
  \node[caixa] (api) {Api};
  \node[caixa, right=1.6cm of api] (app) {Application};
  \node[caixa, right=1.6cm of app, fill=verde!12, draw=verde] (dom) {Domain};
  \node[caixa, below=1.2cm of app] (inf) {Infrastructure};
  \draw[seta] (api) -- (app);
  \draw[seta] (app) -- (dom);
  \draw[seta] (api) |- (inf);
  \draw[seta] (inf) -- (app);
  \draw[seta] (inf) -| (dom);
\end{tikzpicture}
\caption{Regra de dependência entre os projetos.}
\end{figure}

\begin{tabularx}{\textwidth}{@{}lLl@{}}
\toprule
\textbf{Camada} & \textbf{Contém} & \textbf{Depende de} \\
\midrule
Domain         & Entidades com as regras, value objects (\cod{Cpf}), serviços de domínio
                 (\cod{ClassificadorManchester}, \cod{AnalisadorPrescricao}) e interfaces dos repositórios & nada \\
Application    & Casos de uso (\cod{*Service}), DTOs, \cod{ResultWrapper}, chaves de cache,
                 interfaces de infraestrutura & Domain \\
Infrastructure & Repositórios SQL, \cod{DbSession} (Unit of Work), PBKDF2, geração do JWT,
                 \cod{MemoryCacheService} & Application, Domain \\
Api            & Controllers, \cod{Program.cs} (injeção de dependências, JWT, Swagger),
                 middlewares de erro e de cache & Application, Infrastructure \\
\bottomrule
\end{tabularx}

\medskip
O arquivo \cod{HospitalSystem.Domain.csproj} não referencia nenhum outro projeto nem pacote,
o que comprova a regra de dependência.

\subsection{Organização das pastas}

\begin{lstlisting}
HospitalSystem/
+-- api/                              API ASP.NET Core
|   +-- HospitalSystem.sln
|   +-- Dockerfile
|   +-- src/
|       +-- HospitalSystem.Domain/          entidades, regras, interfaces
|       +-- HospitalSystem.Application/     serviços, DTOs, ResultWrapper
|       +-- HospitalSystem.Infrastructure/  SQL, JWT, hash, cache
|       +-- HospitalSystem.Api/             controllers, Program.cs
+-- web/                              aplicativo PHP (cliente da API)
|   +-- lib/ApiClient.php
|   +-- Dockerfile
+-- database/                         scripts SQL
+-- docs/                             regras de negócio e decisões
+-- docker-compose.yml
+-- .env.example
+-- README.md
\end{lstlisting}

\subsection{Fluxo de uma requisição}

Exemplo: \cod{POST /api/internacoes}.

\begin{enumerate}[itemsep=2pt]
  \item \cod{InternacoesController} recebe o JSON e chama o serviço; o controller não tem regra de negócio;
  \item \cod{InternacaoService} busca paciente e leito pelos repositórios e verifica se o paciente já está internado;
  \item a entidade \cod{Internacao} aplica as regras de domínio: o leito precisa estar disponível e passa a ocupado;
  \item o Unit of Work grava internação e leito na mesma transação;
  \item o cache do mapa de leitos e do painel é invalidado;
  \item o serviço devolve um \cod{ResultWrapper}, que o controller converte em \cod{201 Created}.
\end{enumerate}

% ======================= 5 =======================
\section{Execução do projeto}

\subsection{Com Docker (forma oficial)}

Pré-requisito: Docker Desktop instalado e em execução. Na pasta do projeto:

\begin{lstlisting}
copy .env.example .env        (Linux/macOS: cp .env.example .env)
docker compose up --build
\end{lstlisting}

\begin{tabularx}{\textwidth}{@{}lL@{}}
\toprule
\textbf{Serviço} & \textbf{Endereço} \\
\midrule
Aplicativo PHP      & \url{http://localhost:8080/login.php} \\
API (Swagger)       & \url{http://localhost:5080/swagger} \\
Banco MariaDB       & \cod{localhost:3307} (usuário \cod{hospital}, senha do \cod{.env}) \\
\bottomrule
\end{tabularx}

\medskip
Na primeira execução a API cria as tabelas, o administrador e os dados de demonstração. Para
recomeçar do zero: \cod{docker compose down -v}.

\subsection{Execução da API sem Docker}

Pré-requisitos: .NET 10 SDK e MySQL/MariaDB (por exemplo, do XAMPP). As configurações
sensíveis são definidas com \textit{user-secrets}, fora do código:

\begin{lstlisting}
cd api/src/HospitalSystem.Api
dotnet user-secrets set "ConnectionStrings:HospitalDb" "Server=localhost;Port=3306;Database=hospital_api;User ID=root;Password=;"
dotnet user-secrets set "Jwt:Key" "chave-aleatoria-com-pelo-menos-32-caracteres"
dotnet user-secrets set "Seed:AdminSenha" "Admin@2026"
dotnet user-secrets set "Seed:SenhaUsuariosDemo" "Demo@2026"
dotnet run --launch-profile http
\end{lstlisting}

A API sobe em \url{http://localhost:5080} e abre o Swagger.

\subsection{Execução do aplicativo sem Docker}

Com a API rodando, em outro terminal:

\begin{lstlisting}
cd web
php -S localhost:8000
\end{lstlisting}

Acesse \url{http://localhost:8000/login.php}. O endereço da API fica em
\cod{web/config/config.php} (\cod{API\_BASE\_URL}).

\subsection{Usuários de demonstração}\label{sec:usuarios}

\begin{tabularx}{\textwidth}{@{}llL@{}}
\toprule
\textbf{Perfil} & \textbf{E-mail} & \textbf{Senha} \\
\midrule
Administrador & admin@hospital.local      & \cod{Admin@2026} \\
Médico        & medico@hospital.local     & \cod{Demo@2026} \\
Enfermagem    & enfermagem@hospital.local & \cod{Demo@2026} \\
Recepção      & recepcao@hospital.local   & \cod{Demo@2026} \\
\bottomrule
\end{tabularx}

% ======================= 6 =======================
\section{Configuração do banco de dados}

\begin{tabularx}{\textwidth}{@{}lL@{}}
\toprule
\textbf{Configuração} & \textbf{Descrição} \\
\midrule
\cod{ConnectionStrings:HospitalDb}        & conexão com o banco (\cod{.env} no Docker ou \textit{user-secrets}) \\
\cod{Database:InicializarAutomaticamente} & cria banco e tabelas e aplica as migrações ao iniciar \\
\cod{Database:CriarDadosDemo}             & insere dados de exemplo se o banco estiver vazio \\
\bottomrule
\end{tabularx}

\medskip
Scripts na pasta \cod{database/}: \cod{hospital.sql} (schema completo), \cod{dados\_demo.sql}
(dados de demonstração) e \cod{migracao\_php\_para\_api.sql} (migração de um banco do sistema
anterior). Em relação ao banco original, foram acrescentadas a tabela \cod{usuarios} e colunas
em \cod{triagem}, \cod{prontuario}, \cod{prescricoes}, \cod{faturamento} e \cod{uso\_insumos}.

O acesso ao banco usa ADO.NET com SQL parametrizado (\cod{@cpf}, \cod{@nome}), sem concatenar
valores, o que evita SQL Injection. A classe \cod{DbSession} abre uma conexão por requisição e
controla as transações (Unit of Work).

Para consultar o banco no Docker:
\begin{lstlisting}
docker exec -it hospital-db mariadb -u hospital -p hospital
\end{lstlisting}

% ======================= 7 =======================
\section{Autenticação JWT}

\subsection{Funcionamento}

\begin{enumerate}[itemsep=2pt]
  \item O cliente envia \cod{POST /api/auth/login} com e-mail e senha (endpoint público).
  \item A API confere a senha com o hash PBKDF2; a senha nunca é armazenada em texto puro.
  \item Se estiver correta, gera um JWT assinado com HMAC-SHA256 contendo \cod{sub} (id),
        \cod{name}, \cod{email}, \cod{role} (perfil) e \cod{jti}, válido por 60 minutos.
  \item O cliente envia o token em todas as requisições: \cod{Authorization: Bearer <token>}.
  \item A API valida assinatura, emissor, audiência e validade. Sem token ou com token inválido
        a resposta é \textbf{401}; perfil sem permissão recebe \textbf{403}.
\end{enumerate}

\subsection{Ciclo de vida do token}

Emitido no login, enviado em cada requisição, validado pela API sem consultar o banco
(\textit{stateless}), expira após 60 minutos e exige novo login. No logout o aplicativo descarta o
token da sessão.

\subsection{Proteção dos endpoints}

\begin{itemize}[itemsep=2pt]
  \item \textbf{Seguro por padrão}: uma \cod{FallbackPolicy} exige autenticação em todos os
        endpoints, exceto os marcados com \cod{[AllowAnonymous]} (\cod{/api/auth/login} e \cod{/api/status});
  \item autorização por perfil com \cod{[Authorize(Roles = ...)]};
  \item bloqueio da conta por 15 minutos após 5 senhas erradas;
  \item mensagem genérica ``E-mail ou senha inválidos'', para não revelar quais e-mails existem.
\end{itemize}

\subsection{Chaves e configurações sensíveis}

A chave de assinatura (\cod{Jwt:Key}), a string de conexão e as senhas iniciais ficam
\textbf{fora do código-fonte}: no arquivo \cod{.env} (Docker) ou em \textit{user-secrets}
(execução local). O \cod{appsettings.json} versionado mantém esses campos vazios, e o
\cod{.env} está no \cod{.gitignore}. A API não inicia se a chave tiver menos de 32 caracteres.

\subsection{No aplicativo PHP}

\cod{login.php} chama a API e guarda o token em \cod{\$\_SESSION['token']}, no servidor e nunca
no navegador. A classe \cod{ApiClient} envia o token no cabeçalho \cod{Authorization}. Se a API
responder 401, a sessão é encerrada e o usuário volta ao login.

% ======================= 8 =======================
\section{Estratégia de cache}

O cache usa \textbf{IMemoryCache} com a estratégia \textbf{cache-aside}, por trás da interface
\cod{ICacheService} (Application) e da implementação \cod{MemoryCacheService} (Infrastructure):

\begin{enumerate}[itemsep=2pt]
  \item o serviço procura a chave no cache;
  \item \textbf{HIT}: devolve o dado da memória, sem acessar o banco;
  \item \textbf{MISS}: busca no banco, guarda com expiração absoluta e devolve;
  \item toda operação de escrita \textbf{invalida} as chaves afetadas.
\end{enumerate}

\begin{tabularx}{\textwidth}{@{}lllL@{}}
\toprule
\textbf{Dado} & \textbf{Chave} & \textbf{Expiração} & \textbf{Invalidação} \\
\midrule
Médicos, convênios, salas, medicamentos & \cod{catalogo:*}        & 10 min & ao criar, editar ou excluir o item \\
Mapa de leitos                          & \cod{leitos:\{status\}} & 1 min  & mudança em leito ou internação \\
Painel (dashboard)                      & \cod{dashboard}         & 30 s   & mudança em paciente, consulta, triagem, leito ou internação \\
\bottomrule
\end{tabularx}

\medskip
Fila de triagem, estoque e faturamento \textbf{não} usam cache, porque mudam a todo momento e
precisam estar exatos.

\subsection{Como identificar HIT, MISS e invalidação}

Cada resposta traz o cabeçalho \cod{X-Cache: HIT} ou \cod{X-Cache: MISS}, e a API registra no
log:

\begin{lstlisting}
[CACHE MISS] catalogo:medicos → buscando na fonte original (banco de dados)
[CACHE SET]  catalogo:medicos armazenado por 600s
[CACHE HIT]  catalogo:medicos → dado servido da memória (expira em 599s)
[CACHE INVALIDADO] catalogo:medicos → a próxima leitura buscará dados atualizados no banco
[CACHE EXPIROU] dashboard
\end{lstlisting}

No Docker, os logs são vistos com \cod{docker compose logs -f api}. O administrador pode limpar
todo o cache com \cod{DELETE /api/cache}.

% ======================= 9 =======================
\section{Padrão de resposta: ResultWrapper}

Todas as respostas, de sucesso ou de erro, seguem o mesmo formato:

\begin{lstlisting}
{ "success": true, "message": "Operação realizada com sucesso.", "data": { }, "errors": [] }
\end{lstlisting}

\begin{tabularx}{\textwidth}{@{}lL@{}}
\toprule
\textbf{Status} & \textbf{Quando ocorre} \\
\midrule
200 / 201 & sucesso (201 em criação) \\
400 & JSON mal formado ou tipo inválido \\
401 & sem token, token inválido ou expirado, login inválido \\
403 & perfil sem permissão \\
404 & registro não encontrado \\
409 & duplicidade (CPF, CRM, e-mail) \\
422 & regra de negócio violada (mensagem com código \cod{RN-XXX-NN}) \\
500 & erro inesperado (detalhe apenas no log) \\
\bottomrule
\end{tabularx}

\medskip
Erros não tratados são capturados pelo \cod{ExceptionMiddleware} e também saem nesse formato,
sem expor detalhes internos.

% ======================= 10 =======================
\section{Principais endpoints}

A lista completa pode ser testada no Swagger (\url{http://localhost:5080/swagger}). Os endpoints
marcados como públicos não exigem token; todos os demais exigem.

\begin{longtable}{@{}p{2.5cm}p{8.2cm}p{4.4cm}@{}}
\toprule
\textbf{Método} & \textbf{Rota} & \textbf{Descrição} \\
\midrule
\endhead
GET  & \cod{/api/status} & API no ar (público) \\
POST & \cod{/api/auth/login} & autenticação, devolve o JWT (público) \\
GET  & \cod{/api/auth/me} & dados do usuário do token \\
CRUD & \cod{/api/usuarios} & usuários (Administrador) \\
CRUD & \cod{/api/pacientes} & pacientes, busca por nome ou CPF \\
GET  & \cod{/api/pacientes/\{id\}/prontuario} & prontuário completo \\
GET  & \cod{/api/pacientes/\{id\}/gastos} & consolidação de gastos \\
CRUD & \cod{/api/medicos}, \cod{/api/convenios}, \cod{/api/salas}, \cod{/api/medicamentos} & cadastros com cache \\
GET/POST & \cod{/api/consultas} & listar e agendar \\
GET  & \cod{/api/consultas/fila-atendimento} & fila ordenada por Manchester \\
POST & \cod{/api/consultas/\{id\}/triagem} & triagem com risco calculado \\
POST & \cod{/api/prontuarios/\{id\}/prescricoes} & prescrição (Médico) \\
CRUD & \cod{/api/leitos} & leitos (mapa com cache) \\
GET, POST, PATCH & \cod{/api/internacoes} & internação, alta e transferência \\
GET/POST & \cod{/api/insumos} & estoque e movimentações \\
GET, POST, PATCH & \cod{/api/faturamentos} & guias, pagamento, auditoria \\
GET  & \cod{/api/dashboard} & indicadores (com cache) \\
DELETE & \cod{/api/cache} & limpa o cache (Administrador) \\
\bottomrule
\end{longtable}

% ======================= 11 =======================
\section{Exemplos de requisições e respostas}

\subsection{Login}

\begin{lstlisting}
POST /api/auth/login
Content-Type: application/json

{ "email": "admin@hospital.local", "senha": "Admin@2026" }
\end{lstlisting}

Resposta \textbf{200}:
\begin{lstlisting}
{
  "success": true,
  "message": "Login realizado com sucesso.",
  "data": {
    "token": "eyJhbGciOiJIUzI1NiIsInR5cCI6IkpXVCJ9...",
    "tipo": "Bearer",
    "expiraEm": "2026-11-25T14:46:32-03:00",
    "usuario": { "id": 1, "nome": "Administrador do Sistema", "perfil": "Administrador" }
  },
  "errors": []
}
\end{lstlisting}

\subsection{Endpoint protegido sem token}

Resposta \textbf{401}:
\begin{lstlisting}
{ "success": false,
  "message": "Token não informado. Faça login para acessar este recurso.",
  "data": null,
  "errors": ["Token não informado. Faça login para acessar este recurso."] }
\end{lstlisting}

\subsection{Regra de negócio violada}

\cod{POST /api/pacientes} com dados inválidos, resposta \textbf{422}:
\begin{lstlisting}
{ "success": false,
  "message": "A operação viola regras de negócio.",
  "data": null,
  "errors": [
    "RN-PAC-02: CPF inválido.",
    "RN-PAC-04: Data de nascimento é obrigatória e não pode ser futura."
  ] }
\end{lstlisting}

\subsection{Triagem com risco calculado pela API}

\cod{POST /api/consultas/1/triagem}:
\begin{lstlisting}
{ "temperatura": 39.5, "pressaoArterial": "13/9", "frequenciaCardiaca": 105,
  "saturacaoOxigenio": 96, "escalaDor": 3 }
\end{lstlisting}

Resposta \textbf{201}:
\begin{lstlisting}
{ "success": true,
  "message": "Triagem registrada. Classificação: Laranja.",
  "data": { "id": 1, "idConsulta": 1, "classificacaoRisco": "Laranja" },
  "errors": [] }
\end{lstlisting}

\subsection{Perfil sem permissão}

Usuário da Recepção tentando excluir um paciente, resposta \textbf{403}:
\begin{lstlisting}
{ "success": false, "message": "Acesso negado.", "data": null,
  "errors": ["O perfil 'Recepcao' não tem permissão para esta operação."] }
\end{lstlisting}

\subsection{Cache}

\begin{lstlisting}
GET /api/medicos   → 200, X-Cache: MISS
GET /api/medicos   → 200, X-Cache: HIT
POST /api/medicos  → invalida catalogo:medicos
GET /api/medicos   → 200, X-Cache: MISS
\end{lstlisting}

% ======================= 12 =======================
\section{Integração com o aplicativo}

O aplicativo PHP da disciplina anterior deixou de acessar o banco e passou a consumir a API.
A classe \cod{web/lib/ApiClient.php} faz as chamadas HTTP com cURL, envia o token JWT no
cabeçalho \cod{Authorization} e lê o \cod{ResultWrapper}: se \cod{success} for verdadeiro, usa
\cod{data}; caso contrário, mostra \cod{message} e \cod{errors} na tela. Se a API estiver fora do
ar, o aplicativo exibe uma mensagem amigável.

% ======================= 13 =======================
\section{Principais decisões arquiteturais}

\begin{tabularx}{\textwidth}{@{}lL@{}}
\toprule
\textbf{Decisão} & \textbf{Motivo} \\
\midrule
Clean Architecture em 4 projetos & isola as regras de negócio da tecnologia; trocar banco ou cache mexe só na Infrastructure \\
Regras dentro das entidades & propriedades com \cod{private set}; não é possível criar estado inválido sem passar pela regra \\
ADO.NET sem Entity Framework & o banco já existia; com SQL explícito o schema original foi mantido \\
Unit of Work & operações que precisam acontecer juntas (internar e ocupar leito) ficam numa transação \\
ResultWrapper único & o cliente trata todas as respostas do mesmo jeito \\
JWT + PBKDF2 & autenticação sem estado no servidor e senhas resistentes a força bruta \\
IMemoryCache atrás de \cod{ICacheService} & simples para uma instância; pode virar Redis sem mudar os serviços \\
Segredos fora do código & \cod{.env} e \textit{user-secrets}; nada sensível vai para o Git \\
Docker Compose & banco, API e aplicativo sobem com um único comando \\
\bottomrule
\end{tabularx}

\subsection{Escalabilidade}

A API não guarda sessão no servidor (o JWT carrega a identidade do usuário) e roda em
container, então é possível executar várias instâncias atrás de um balanceador de carga. Para
isso, dois pontos mudariam: o cache em memória passaria a ser compartilhado (Redis, trocando
apenas a implementação de \cod{ICacheService}) e a criação do banco sairia da inicialização da
API para uma etapa separada de migração.

% ======================= 14 =======================
\section{Reprodução do projeto}

\begin{lstlisting}
git clone <URL do repositório>
cd HospitalSystem
copy .env.example .env
docker compose up --build
\end{lstlisting}

Depois, acessar \url{http://localhost:5080/swagger} (API) e
\url{http://localhost:8080/login.php} (aplicativo) com os usuários da seção~\ref{sec:usuarios}.

\end{document}
